% =========================================================
% Computational Linear Algebra --- Determinants and the Cross Product
% (Strang \S{}5.2 / Johnston \S{}1.A / Ashraf \S{}1.4)
% =========================================================

\subsection{Determinants and the Cross Product}

% ---------------------------------------------------------
\subsection*{\texorpdfstring{The $3 \times 3$ Determinant}{The 3 x 3 Determinant}}

\begin{definitionbox}{Definition ($3 \times 3$ Determinant via Cofactor Expansion)}
\begin{definition}[$3 \times 3$ Determinant via Cofactor Expansion]
\label{def:three-by-three-determinant-cofactor-expansion}
For $A \in \mathbb{R}^{3 \times 3}$, expanding along the first row:
The signs alternate as $+, -, +$ along the first row.
\[
  \det(A) = a_{11}(a_{22}a_{33} - a_{23}a_{32})
           - a_{12}(a_{21}a_{33} - a_{23}a_{31})
           + a_{13}(a_{21}a_{32} - a_{22}a_{31}).
\]
\end{definition}
\end{definitionbox}
\begin{remark*}[Standard quantified statement]
This definition fixes the named object, operation, or relation by the formula
or conditions stated in the definition.
\end{remark*}

\begin{remark*}[Interpretation]
The statement records the local meaning of the named concept used by the
surrounding computational or geometric construction.
\end{remark*}

\DefinitionalRoot

% ---------------------------------------------------------
\subsection*{The Cross Product}

\begin{definitionbox}{Definition (Cross Product)}
\begin{definition}[Cross Product]
\label{def:cross-product}
For $\mathbf{v} = (v_1, v_2, v_3)$ and $\mathbf{w} = (w_1, w_2, w_3)$ in $\mathbb{R}^3$,
the \emph{cross product} is
\[
  \mathbf{v} \times \mathbf{w}
  = \begin{vmatrix} \mathbf{e}_1 & \mathbf{e}_2 & \mathbf{e}_3 \\
                    v_1 & v_2 & v_3 \\
                    w_1 & w_2 & w_3
    \end{vmatrix}
  = (v_2 w_3 - v_3 w_2,\;\; v_3 w_1 - v_1 w_3,\;\; v_1 w_2 - v_2 w_1).
\]
\end{definition}
\end{definitionbox}
\begin{remark*}[Standard quantified statement]
This definition fixes the named object, operation, or relation by the formula
or conditions stated in the definition.
\end{remark*}

\begin{remark*}[Interpretation]
The statement records the local meaning of the named concept used by the
surrounding computational or geometric construction.
\end{remark*}

\DefinitionalRoot

\begin{remark*}[Exposition]
\textbf{Properties of the Cross Product.}
For $\mathbf{v}, \mathbf{w}, \mathbf{x} \in \mathbb{R}^3$ and $c \in \mathbb{R}$:
\begin{enumerate}
  \item $\mathbf{v} \times \mathbf{w} = -(\mathbf{w} \times \mathbf{v})$ \hfill (anticommutativity)
  \item $\mathbf{v} \times (\mathbf{w} + \mathbf{x}) = \mathbf{v} \times \mathbf{w} + \mathbf{v} \times \mathbf{x}$ \hfill (distributivity)
  \item $\mathbf{v} \times \mathbf{v} = \mathbf{0}$
  \item $(c\mathbf{v}) \times \mathbf{w} = c(\mathbf{v} \times \mathbf{w})$
  \item $\mathbf{v} \cdot (\mathbf{v} \times \mathbf{w}) = 0$ and $\mathbf{w} \cdot (\mathbf{v} \times \mathbf{w}) = 0$ \hfill (orthogonality)
  \item $\|\mathbf{v} \times \mathbf{w}\| = \|\mathbf{v}\|\,\|\mathbf{w}\|\sin\theta$ \hfill (magnitude = parallelogram area)
\end{enumerate}
\end{remark*}

\begin{remark*}[Exposition]
\textbf{NURBS connection: surface normal vectors.}
The unit normal to a parametric surface at a point $(u,v)$ is
\[
  N(u,v) = \frac{S_u \times S_v}{\|S_u \times S_v\|},
\]
where $S_u$ and $S_v$ are the partial derivative vectors of the surface
with respect to the parameters.
The cross product gives the direction orthogonal to both tangent directions;
dividing by the norm produces the unit normal.
This appears in Piegl \& Tiller \S{}1.1 and is used throughout whenever
normal vectors are needed (shading, offset surfaces, etc.).
\end{remark*}

\begin{remark*}[Exposition]
\textbf{Mnemonic for the cross product.}
Write the entries of $\mathbf{v}$ and $\mathbf{w}$ each twice in a $2 \times 6$ array.
Ignore the first and last columns.
Draw three ``X'' shapes between columns 2--3, 3--4, and 4--5.
Each ``X'' contributes one entry of $\mathbf{v} \times \mathbf{w}$: add along forward
diagonals, subtract along backward diagonals.
\end{remark*}
